\subsection{DEFINITION OF PRIMITIVES}

A vector function \(\mathbf{f}\) defined on an interval \(\mathrm{I} \subset \mathbf{R}\) cannot be the derivative at every point of this interval of a vector function \(\mathbf{g}\) (defined and continuous on \(\mathrm{I}\) ) unless it satisfies quite stringent conditions: for example, if \(\mathbf{f}\) admits a right limit and a left limit at a point \(x_0\) interior to \(\mathrm{I}\) then \(\mathbf{f}\) must be continuous at the point \(x_0\) , as follows from prop. 6 of \(\mathrm{I}\) , p.~18; it follows, that if one takes the interval \([-1, 1]\) for \(\mathrm{I}\) , and for \(f\) the real function equal to \(-1\) on \([-1, 0[\) , and to \(+1\) on \([0, 1]\) , then \(f\) is not the derivative of any continuous function on \(\mathrm{I}\) ; all the same, the function \(|x|\) has \(f(x)\) as its derivative at every point \(\neq 0\) ; one is thus led to make the following definition:

DEFINITION 1. Given a vector function f defined on an interval I ⊂ R we say that a function g defined on I is a primitive of f if g is continuous on I and has a derivative equal to f(x) at every point x of the complement (with respect to I) of a countable subset of I.

If also \(\mathbf{g}\) admits a derivative equal to \(\mathbf{f}(x)\) at every point \(x\) of I, one says that \(\mathbf{g}\) is a strict primitive of \(\mathbf{f}\) .

With this definition, one sees that the real function \(f\) considered above admits a primitive equal to \(|x|\) .

It is clear that if f admits a primitive on I then every primitive of f is also a primitive of every function which is equal to f except at the points of a countable subset of I. By an abuse of language one speaks of a primitive on I of a function \(f_{0}\) defined only on the complement (with respect to I) of a countable subset of I: this will be the primitive of every function f defined on I and equal to \(f_{0}\) at the points where \(f_{0}\) is defined.

PROPOSITION 1. Let f be a vector function defined on I with values in E; if f admits a primitive g on I then the set of primitives of f on I is identical to the set of functions \(g + a\) , where a is a constant function with its values in E.

Indeed, it is clear that \(\mathbf{g} + \mathbf{a}\) is a primitive of \(\mathbf{f}\) for any \(\mathbf{a} \in E\) ; on the other hand, if \(\mathbf{g}_1\) is a primitive of \(\mathbf{f}\) then \(\mathbf{g}_1 - \mathbf{g}\) admits a derivative equal to 0 except at the points of a countable subset of I, and thus is constant (I, p.~17, corollary).

One says that the primitives of a function f (when they exist) are defined ``up to an additive constant''. To define a primitive of f unambiguously it is enough to assign it (arbitrarily) a value at a point \(x_{0} \in I\) ; in particular, there exists one and only one primitive g of f such that \(\mathbf{g}(x_{0}) = 0\) ; for every primitive h of f one has \(\mathbf{g}(x) = \mathbf{h}(x) - \mathbf{h}(x_{0})\) .

